\documentclass[12pt]{article}
\usepackage[T1]{fontenc}
\usepackage[utf8]{inputenc}
\usepackage{lmodern}
\usepackage{textcomp}
\usepackage{enumitem}
\usepackage{geometry}
\geometry{margin=1in}
\begin{document}
\begin{center}
Answer Key - Biology - Undergraduate Level Assessment 19 \
\small (Answers are ordered exactly as in the question paper.)
\end{center}

\section*{Multiple Choice}
\begin{enumerate}[label=\arabic*.]
\item \textbf{Answer:} D
\\ \textit{Options:} A) Earth is millions of years old, and populations are unchanging.; B) Earth is millions of years old, and populations gradually change.; C) Earth is a few thousand years old, and populations rapidly change.; D) Earth is a few thousand years old, and populations are unchanging.; E) Earth is billions of years old, and populations gradually change.
\\ \textit{Note:} Prior to the time of Lyell and Darwin, the prevailing notion was based on a literal interpretation of biblical texts, which suggested that the Earth was created only a few thousand years ago and that species were immutable and unchanged over time, contrasting sharply with the later concepts of deep geological time and evolutionary change.
\item \textbf{Answer:} D
\\ \textit{Options:} A) Epidermal cell; B) Stomatal guard cell; C) Root cap cell; D) Xylem vessel member; E) Phloem sieve tube member
\\ \textit{Note:} Xylem vessel members undergo programmed cell death to lose their cellular components, allowing for the formation of hollow tubes that efficiently transport water and nutrients in plants.
\item \textbf{Answer:} D
\\ \textit{Options:} A) O$_{2}$; B) CH 4 (Methane); C) CO 2; D) H$_{2}$O; E) N$_{2}$ (Nitrogen)
\\ \textit{Note:} The correct answer is H$_{2}$O because during photosynthesis, water molecules are split to release oxygen as a byproduct while the remaining hydrogen is used in the formation of glucose.
\item \textbf{Answer:} A
\\ \textit{Options:} A) Ferns grow larger because they can perform photosynthesis more efficiently than mosses; B) Ferns grow larger due to higher nutrient content; C) Ferns grow larger due to a stronger stem; D) Ferns grow larger due to their ability to move to more fertile areas; E) Mosses are limited in size by their shorter life cycle compared to ferns
\\ \textit{Note:} Ferns possess a well-developed vascular system that allows for more efficient transport of water and nutrients, enabling them to perform photosynthesis at a greater scale than mosses, which lack this system and are limited in size.
\item \textbf{Answer:} A
\\ \textit{Options:} A) oxygen; B) calcium; C) carbon; D) potassium; E) phosphorus
\\ \textit{Note:} Oxygen is the principle inorganic compound found in living organisms because it is essential for cellular respiration and is a major component of water, which is critical for life processes.
\end{enumerate}

\section*{True / False}
\begin{enumerate}[label=\arabic*.]
\setcounter{enumi}{5}
\item \textbf{Answer:} False
\\ \textit{Options:} A) True; B) False
\\ \textit{Note:} High relative humidity decreases the rate of transpiration from plants because it reduces the water vapor gradient between the leaf surfaces and the surrounding air, making it more difficult for water to evaporate.
\item \textbf{Answer:} False
\\ \textit{Options:} A) True; B) False
\\ \textit{Note:} An action potential is characterized by a rapid depolarization followed by repolarization, typically moving from a resting membrane potential of around -70 millivolts to a peak of about +30 millivolts, rather than a direct change from +50 millivolts to -70 millivolts.
\item \textbf{Answer:} True
\\ \textit{Options:} A) True; B) False
\\ \textit{Note:} The statement is true because molting in crustaceans is regulated by hormones such as ecdysteroids and is triggered by environmental cues like temperature, light, and food availability, which influence their growth and development.
\item \textbf{Answer:} False
\\ \textit{Options:} A) True; B) False
\\ \textit{Note:} The epidermis in plants primarily serves as a protective barrier and does not facilitate gas exchange and nutrient transport; instead, gas exchange occurs through specialized structures called stomata, and nutrient transport is primarily conducted by vascular tissues.
\item \textbf{Answer:} False
\\ \textit{Options:} A) True; B) False
\\ \textit{Note:} The statement is false because schizophrenia has a significant genetic component, with studies indicating that heritability accounts for a substantial portion of the risk for developing the disorder, alongside environmental factors.
\end{enumerate}

\section*{Long Form Response}
\begin{enumerate}[label=\arabic*.]
\setcounter{enumi}{10}
\item \textbf{Answer:} Inbreeding and loss of genetic variation significantly increase the vulnerability of small populations, often leading them toward an extinction vortex. Inbreeding occurs when closely related individuals mate, resulting in offspring with a higher probability of inheriting deleterious alleles, which can diminish overall fitness. Additionally, a lack of genetic variation restricts the population's ability to adapt to environmental changes, making them more susceptible to diseases and changing conditions. This combination can create a feedback loop, where declining population size leads to further inbreeding and reduced genetic diversity, ultimately pushing the population closer to extinction. Thus, maintaining genetic diversity is crucial for the survival and resilience of small populations.
\\ \textit{Note:} correct because it effectively connects inbreeding and reduced genetic variation to decreased fitness and adaptability in small populations, illustrating how these factors can lead to a cycle of decline and potential extinction.
\item \textbf{Answer:} Combustion and aerobic respiration are two distinct processes that release energy, but they differ significantly in the forms of energy produced and their implications for biological systems. Combustion primarily releases energy as light and heat, resulting in rapid energy output, which is useful for applications like engines and fires but can be destructive and unsustainable in biological contexts. In contrast, aerobic respiration, which occurs in living organisms, releases energy predominantly as heat through the breakdown of glucose in the presence of oxygen. This method of energy release is more controlled and efficient, allowing organisms to harness energy gradually for cellular processes, thus supporting life functions without the destructive consequences associated with combustion. The heat produced in aerobic respiration also plays a critical role in maintaining body temperature in warm-blooded animals, highlighting its essential role in biological systems.
\\ \textit{Note:} correct because it clearly distinguishes between combustion, which produces energy as light and heat in a rapid and often destructive manner, and aerobic respiration, which efficiently generates energy for biological processes in a controlled manner, highlighting their differing implications for living systems.
\end{enumerate}

\vfill
\noindent\textit{End of Answer Key}
\end{document}